%% Code listings and algorithm pseudocode
%% Two workflows:
%%   1. `code` / `codebox`  — concrete syntax (minted + tcolorbox)
%%   2. `algorithm`         — pseudocode (algpseudocodex + tcolorbox)

%%% ────────────────────────────────────────────────
%%% Minted — syntax-highlighted code listings
%%% ────────────────────────────────────────────────

\usepackage{minted}
\tcbuselibrary{minted}

%% Global minted defaults
\setminted{
    fontsize=\small,
    baselinestretch=1.1,
    breaklines,              % allow automatic line breaks
    breakanywhere=false,     % break only at whitespace
    tabsize=4,
    numbersep=8pt,
    xleftmargin=0.6em,
    xrightmargin=0.6em,
}

%% Inline code with syntax highlighting: \mintinline or the short \code macro
\NewDocumentCommand{\code}{O{text} v}{\mintinline{#1}{#2}}

%%% tcolorbox wrapper for minted — "codebox"
%%% Provides a framed, breakable box that matches the look of theorem-like environments.
\tcbset{
    codebox/.style={
        breakable,
        parbox=false,
        colframe=black!50,
        colback=LightGrayFill,
        boxrule=0.4pt,
        arc=3pt,
        left=0.5em, right=0.5em,
        top=0.8ex, bottom=0.8ex,
        beforeafter skip=0.5\baselineskip plus 2pt,
    },
}

%% Floating code listing environment
%% Usage:
%%   \begin{codebox}[float,caption={My listing},label={lst:example}]{julia}
%%     function foo(x)
%%         return x^2
%%     end
%%   \end{codebox}
%%
%% Non-floating (inline) usage:
%%   \begin{codebox}{julia}
%%     x = 1 + 2
%%   \end{codebox}
%%
%% With line numbers:
%%   \begin{codebox}[linenos]{julia}
%%     ...
%%   \end{codebox}
\DeclareTCBListing{codebox}{ O{} m }{%
    codebox,
    listing engine=minted,
    minted language=#2,
    minted options={fontsize=\small, baselinestretch=1.1, breaklines},
    listing only,
    #1, % user options (float, caption, label, linenos, ...)
}

%% Share the `table` counter for code listings (consistent with figures/tables/equations/theorems)
\makeatletter
\AtBeginDocument{%
    \@ifundefined{c@tcb@cnt@codebox}{}{%
        \let\c@tcb@cnt@codebox\c@table
    }%
}
\makeatother

%% zref-clever type name for code listings
\zcRefTypeSetup{tcb@cnt@codebox}{
    name-sg = listing, Name-sg = Listing,
    name-pl = listings, Name-pl = Listings,
}


%%% ────────────────────────────────────────────────
%%% Pseudocode — algpseudocodex
%%% ────────────────────────────────────────────────

\usepackage{algorithm}           % float wrapper for algorithms
\usepackage[noEnd=false]{algpseudocodex} % pseudocode typesetting with indent guides

%% Share the `table` counter for algorithm floats
\makeatletter
\let\c@algorithm\c@table
\makeatother

%% zref-clever type name for algorithms
\zcRefTypeSetup{algorithm}{
    name-sg = algorithm, Name-sg = Algorithm,
    name-pl = algorithms, Name-pl = Algorithms,
}

%% Visual style for the algorithm float — match the look of theorem-like tcolorbox environments
%% We wrap the algorithm body in a tcolorbox via \floatsetup
\floatsetup[algorithm]{
    style=plain,
}

%% Style the algorithm caption to match other floats
\captionsetup[algorithm]{
    format=plain,
    font=small,
    labelfont={sf,bf},
    labelsep=slash,
}

%% Sans-serif keyword style to match theorem headings
\algrenewcommand\algorithmicfunction{\textsf{\textbf{function}}}
\algrenewcommand\algorithmicprocedure{\textsf{\textbf{procedure}}}
\algrenewcommand\algorithmicif{\textsf{\textbf{if}}}
\algrenewcommand\algorithmicthen{\textsf{\textbf{then}}}
\algrenewcommand\algorithmicelse{\textsf{\textbf{else}}}
\algrenewcommand\algorithmicfor{\textsf{\textbf{for}}}
\algrenewcommand\algorithmicwhile{\textsf{\textbf{while}}}
\algrenewcommand\algorithmicdo{\textsf{\textbf{do}}}
\algrenewcommand\algorithmicreturn{\textsf{\textbf{return}}}
\algrenewcommand\algorithmicend{\textsf{\textbf{end}}}
\algrenewcommand\algorithmicrequire{\textsf{\textbf{Require:}}}
\algrenewcommand\algorithmicensure{\textsf{\textbf{Ensure:}}}

%% tcolorbox wrapper for the algorithm body (non-float, for use inside algorithm environment)
%% Usage:
%%   \begin{algorithm}
%%     \caption{Euler step}\label{alg:euler}
%%     \begin{algobox}
%%       \begin{algorithmic}[1]
%%         \Procedure{EulerStep}{$u, \Delta t$}
%%           \State $u \gets u + \Delta t \cdot f(u)$
%%         \EndProcedure
%%       \end{algorithmic}
%%     \end{algobox}
%%   \end{algorithm}
\newtcolorbox{algobox}{
    codebox, % reuse the same visual style
    before skip=0pt, after skip=0pt, % no extra space inside the float
}
